\section*{Capítulo \ref{ch:b1:nucleic-acids} --- Nucleotídeos e ácidos nucleicos}

\begin{solution}{exo:b1:nucleic-acids:1}
Uma \omterm{def:b1:nucleic-acids:nucleotide}{base nitrogenada}, uma pentose, de um a três fosfatos. \omterm{def:b1:nucleic-acids:chain}{RNA}: ribose
(com uma hidroxila $2'$) e uracila; \omterm{def:b1:nucleic-acids:chain}{DNA}: $2'$-desoxirribose e timina.
\end{solution}

\begin{solution}{exo:b1:nucleic-acids:2}
$5'$-TAGGATCC-$3'$. Pares: G--C, G--C, A--T, T--A, C--G, C--G, T--A,
A--T: quatro G--C (12 ligações) e quatro A--T (8): 20 \omterm{def:b1:water-small-molecules:water}{ligações de
hidrogênio}.
\end{solution}

\begin{solution}{exo:b1:nucleic-acids:3}
T \qty{22}{\%}; G e C dividem os \qty{56}{\%} restantes: \qty{28}{\%}
cada; $G + C = \qty{56}{\%}$.
\end{solution}

\begin{solution}{exo:b1:nucleic-acids:4}
Entre as duas curvas, a \qty{0.4}{\celsius} por ponto percentual: cerca
de \qty{85}{\celsius}. A absorbância sobe \qty{40}{\%}.
\end{solution}

\begin{solution}{exo:b1:nucleic-acids:5}
$\num{4.6e6}\times 0.34\,\mathrm{nm} = \qty{1.56}{mm}$; \num{460000}
voltas; $1560/2 = 780$ vezes em comprimento (o cromossomo é dobrado em
alças e super-hélices).
\end{solution}

\begin{solution}{exo:b1:nucleic-acids:6}
As \omterm{def:b1:nucleic-acids:nucleotide}{bases} pareiam com as bordas voltadas uma para a outra, e o açúcar de
cada uma fica preso do mesmo lado do par apenas se os dois esqueletos
correrem em sentidos opostos; com fitas paralelas os dois açúcares
ficariam em bordas opostas, e a largura do par, e a posição do
esqueleto, variariam com a sequência. Fitas antiparalelas com pares
purina--pirimidina dão os \qty{2}{nm} constantes que o padrão da fibra
exige.
\end{solution}

\begin{solution}{exo:b1:nucleic-acids:7}
$0.35\times 50 = \qty{17.5}{\micro g/mL}$. $A_{260}/A_{280} = 1.75$:
\omterm{def:b1:nucleic-acids:chain}{DNA} essencialmente puro (1.8), talvez com um traço de proteína. Em
\qty{1}{mL}: $17.5\times 10^{-6}/(650\times 1.66\times 10^{-24}) =
\num{1.6e16}$ \omterm{thm:b1:nucleic-acids:helix}{pares de bases}.
\end{solution}

\begin{solution}{exo:b1:nucleic-acids:8}
As \omterm{def:b1:nucleic-acids:nucleotide}{bases} empilhadas se protegem umas às outras e absorvem menos;
separar as fitas expõe cada \omterm{def:b1:nucleic-acids:nucleotide}{base} à luz e a absorbância sobe
\qty{40}{\%}. Um \omterm{def:b1:nucleic-acids:chain}{DNA} puro funde dentro de poucos graus em torno do
$T_m$ único dele, de modo cooperativo; uma mistura contém \omterm{def:b1:nucleic-acids:chain}{DNAs} de muitos
teores de $G + C$, cada um com o próprio $T_m$, de modo que a soma das
curvas deles se espalha por dezenas de graus.
\end{solution}

\begin{solution}{exo:b1:nucleic-acids:9}
Logo abaixo de \qty{62}{\celsius} e acima de \qty{57}{\celsius},
digamos \qty{60}{\celsius}: o híbrido perfeito se mantém, e o descasado
funde. A \qty{45}{\celsius} os híbridos com até três descasamentos são
estáveis e a sonda também acende sequências aparentadas.
\end{solution}

\begin{solution}{exo:b1:nucleic-acids:10}
As regras decorrem do pareamento de \omterm{def:b1:nucleic-acids:nucleotide}{bases} entre duas fitas
complementares. O \omterm{def:b1:nucleic-acids:chain}{RNA}, de fita simples, não tem parceira que restrinja a
composição dele, de modo que $A \neq U$ em geral; os vírus pequenos em
questão têm genomas de \omterm{def:b1:nucleic-acids:chain}{DNA} de fita simples, que também não precisam
obedecer às regras --- e o descumprimento delas foi uma das primeiras
pistas de que esses genomas eram de fita simples.
\end{solution}

\begin{solution}{exo:b1:nucleic-acids:11}
Cada par sozinho vale alguns quilojoules por mol, menos que a agitação
térmica, e se forma e se rompe em microssegundos. Numa hélice os pares
se formam juntos: cada par se empilha sobre o anterior e os primeiros
pares nucleiam os demais, de modo que as energias se somam e a
probabilidade de todos se romperem de uma vez é desprezível --- é a
cooperatividade. O empilhamento contribui tanto quanto as \omterm{def:b1:water-small-molecules:water}{ligações de
hidrogênio}. Para uma sonda isso significa que a estabilidade cresce com
o comprimento e cai bruscamente a cada descasamento, que interrompe a
sequência cooperativa: a especificidade é uma propriedade do trecho
inteiro.
\end{solution}

\begin{solution}{exo:b1:nucleic-acids:12}
A hélice explica a cópia de imediato: cada fita especifica a outra, de
modo que separá-las e parear \omterm{def:b1:nucleic-acids:nucleotide}{nucleotídeos} novos em cada uma dá duas
moléculas idênticas; uma mutação é uma \omterm{def:b1:nucleic-acids:nucleotide}{base} errada fixada numa cópia e
fielmente copiada dali em diante. Ela não diz, sozinha, como uma
sequência de quatro \omterm{def:b1:nucleic-acids:nucleotide}{bases} especifica uma proteína --- isso levou uma
década (código, mRNA, tRNA, ribossomo). O modelo foi aceito porque se
encaixava em todas as medidas (Chargaff, as dimensões dos raios X, a
densidade) sem sobrar nada, e porque a explicação que ele dava da
hereditariedade era tão econômica que a alternativa --- uma coincidência
--- era implausível; Meselson e Stahl confirmaram depois a cópia que ele
previa.
\end{solution}

\begin{solution}{pb:b1:nucleic-acids:1}
\textbf{1.} $\num{48502}\times 0.34 = \qty{16490}{nm} = \qty{16.5}{\micro m}$.
\textbf{2.} $\num{48502}/10 = 4850$ voltas.
\textbf{3.} $\num{48502}\times 650 = \qty{3.15e7}{Da}$.
\textbf{4.} $\num{3.15e7}\times 1.66\times 10^{-24} = \qty{5.2e-17}{g}$.
\textbf{5.} Dois por par: \num{97004} fosfatos, carga líquida de
$-97\,004\,e$.
\textbf{6.} $G + C = 0.499\times\num{97004} = \num{48405}$: $G = C =
\num{24203}$; $A = T = \num{24299}$.
\textbf{7.} $3\times\num{24203} + 2\times\num{24299} = \num{121207}$.
\textbf{8.} $81.5 + 16.6\times(-1) + 0.41\times 49.9 - 500/48502 =
81.5 - 16.6 + 20.5 - 0.01 = \qty{85.4}{\celsius}$.
\textbf{9.} $81.5 - 33.2 + 20.5 = \qty{68.8}{\celsius}$. Os cátions
blindam a repulsão entre os dois esqueletos carregados negativamente;
com menos sal as fitas se repelem mais e se separam a uma temperatura
menor.
\textbf{10.} \qty{35}{\%}: $81.5 - 16.6 + 14.4 - 0.5 =
\qty{78.8}{\celsius}$; \qty{62}{\%}: $81.5 - 16.6 + 25.4 - 0.5 =
\qty{89.8}{\celsius}$. A molécula inteira funde em etapas: as regiões
ricas em A--T se abrem primeiro, e as ricas em G--C por último, de modo
que a curva é uma subida larga com ombros, e não um degrau único e
brusco.
\textbf{11.} $0.50\times 50 = \qty{25}{\micro g/mL}$; depois da fusão
$A_{260} = 0.70$.
\textbf{12.} $25\times 10^{-6}/5.2\times 10^{-17} = \num{4.8e11}$ moléculas.
\textbf{13.} $81.5 - 16.6 + 0.41\times 50 - 500/12 = 81.5 - 16.6 + 20.5
- 41.7 = \qty{43.7}{\celsius}$ (a fórmula é grosseira para trechos tão
curtos). A \qty{37}{\celsius} as extremidades pareiam, mas se abririam e
fechariam; na \omterm{def:b1:cell-unit-of-life:cell}{célula} uma enzima (a ligase) sela covalentemente as duas
falhas, e o círculo deixa de depender dos doze pares.
\textbf{14.} $\pi\times(1.0)^2\times\num{16490} = \qty{5.2e4}{nm^3}$.
\textbf{15.} $\frac{4}{3}\pi\times 27.5^3 = \qty{8.7e4}{nm^3}$.
\textbf{16.} \qty{60}{\%}: dois terços do empacotamento mais compacto
possível de cilindros --- o \omterm{def:b1:nucleic-acids:chain}{DNA} é enrolado em espirais quase
cristalinas.
\textbf{17.} Contraíons (poliaminas, $\mathrm{Mg^{2+}}$) precisam
neutralizar quase toda a carga, ou a repulsão não poderia ser vencida; a
repulsão restante e a energia de dobramento armazenada nas espirais
apertadas empurram para fora contra a casca, e quando a cauda se abre
essa pressão impele a primeira parte do genoma para dentro da bactéria.
\textbf{18.} As alças num capsídeo de \qty{55}{nm} têm raios de
\qtyrange{10}{25}{nm}, abaixo dos \qty{50}{nm} em que o \omterm{def:b1:nucleic-acids:chain}{DNA} se dobra
livremente: dobrá-lo custa energia, paga pelo motor de empacotamento
(\omterm{def:b1:water-small-molecules:atp}{ATP}) e armazenada como a pressão da questão 17.
\textbf{19.} $5'$-AGGTCGCCGCCC-$3'$.
\textbf{20.} \num{1200} \omterm{def:b1:nucleic-acids:nucleotide}{nucleotídeos}; $1200\times 330 = \qty{4.0e5}{Da}$.
\textbf{21.} $75\times 600 = \num{45000}$ pares, \qty{93}{\%}
codificantes --- um genoma compacto; o genoma humano é sessenta vezes
menos denso em genes.
\textbf{22.} Duas forquilhas, cada uma copiando \num{24251} pares em
\qty{1200}{s}: \qty{20}{} pares por segundo por forquilha (a polimerase
do hospedeiro é capaz de \qty{1000}{}; o tempo é fixado pelo ciclo
inteiro, e não pela polimerase).
\textbf{23.} $\num{121207}/\num{48502} = 2.5$ ligações por par; duas forquilhas a 20 pares por segundo: 100 \omterm{def:b1:water-small-molecules:water}{ligações de hidrogênio} por segundo --- e outras tantas refeitas atrás delas.
\textbf{24.} Um par em \num{48502} muda o $G + C$ em \qty{0.002}{\%} e o
$T_m$ em \qty{0.001}{\celsius}: indetectável fundindo a molécula
inteira; só uma sonda curta que cobrisse o sítio o mostraria.
\textbf{25.} $T_m = \qty{85}{\celsius}$ em sódio a \qty{0.1}{mol/L};
comprimento de contorno de \qty{16.5}{\micro m}, trezentas vezes o
diâmetro do capsídeo, empacotado nele a \qty{60}{\%} do volume dele.
\end{solution}
